\begin{enumerate}

\item La constante de torsión es el momento de inercia polar de la corona.

\item La tensión máxima está en el radio exterior.

\item El giro total es el giro unitario por la longitud.

\end{enumerate}

\noindent\textbf{Operaciones y resultados intermedios.}

\begin{enumerate}

\item Constante de torsión.

\[\fitmath{y_{1}=\frac{\pi\cdot \left({25}^{4}-{20}^{4}\right)}{2}\approx 362264.9\,\mathrm{mm^{4}}}\]

\item Tensión tangencial máxima.

\[\fitmath{y_{2}=\frac{1000000\cdot 25}{\frac{\pi\cdot \left({25}^{4}-{20}^{4}\right)}{2}}\approx 69.01027\,\mathrm{MPa}}\]

\item Giro entre extremos.

\[\fitmath{y_{3}=\frac{1000000\cdot 1000}{\frac{26950\cdot \pi\cdot \left({25}^{4}-{20}^{4}\right)}{2}}\approx 0.1024271\,\mathrm{rad}}\]

\end{enumerate}
